\section{Discussion}
\label{sec:discussion}

The per-configuration breakdown in Table~\ref{tab:iterations} concentrates 80\% of the total improvement (12.98 of 16.15 nDCG points) in the single transition from Approach~2 to Approach~3. That is from a frozen-feature ensemble to an end-to-end fine-tuned ViT-L driven by soft-label supervision. The asymmetry across the configurations is itself the central observation, and it can be unpacked along three connected lines: the low ceiling of inference-time refinements over frozen backbones, the dominant effect of the bundled encoder-and-loss swap, and the diminishing returns of post-fine-tuning refinements that motivate the future-work direction.

Approach~2 fuses three pre-trained models that were never optimized for the soft-label setting: JPoSE and MMEN are triplet-supervised, MI-MM is binary multi-instance MaxMargin, and none sees the relevance matrix $R$ during training. Equal-weight averaging of their normalized scores recovers $+1.78$~nDCG on top of JPoSE, in line with the canonical observation that retrieval ensembles reduce variance without correcting bias~\cite{wortsman2022modelsoups}. Because the three models share the same failure mode (they ignore graded relevance), averaging them produces a more stable estimate of the same flawed objective rather than one that is closer to nDCG. Equal weights were used in the absence of a relevance-aligned validation signal and weighted averaging tuned on a held-out split could plausibly add a fraction of a point but is unlikely to overcome the bias ceiling. Graph-diffusion re-ranking adds another $+0.51$~nDCG at no training cost, consistent with magnitudes reported for diffusion in image retrieval~\cite{iscen2017diffusion}. A small $k=20$ keeps the affinity graph sparse and avoids over-smoothing, moderate $\alpha=0.3$ blends rather than replaces the original signal, and the symmetric construction in Eq.~\eqref{eq:rerank}---propagating through both visual and textual graphs---prevents either modality from dominating the result.

The single largest jump, $+12.98$~nDCG between configurations~3 and~4, is associated with a \emph{bundle of changes} introduced together at that step: pre-extracted features are replaced by an end-to-end ViT-L/14 encoder pre-trained on egocentric video~\cite{zhao2023avion,zhao2023lavila}, binary triplet/MaxMargin objectives are replaced by the SMS loss~\cite{wang2024symmetricmultisimilaritylossepickitchens100} that consumes the soft-label matrix $R$ directly, and both encoders are jointly fine-tuned rather than frozen. Because these changes were applied jointly and not ablated, we report the gain as a property of the bundle rather than attributing it causally to any single factor; disentangling the relative contribution of each ingredient would require an ablation study that was beyond our compute budget. What can be stated empirically is that the bundle dominates every other change in the table by an order of magnitude, and that the only mechanism in this bundle that directly optimises the nDCG metric (Eq.~\eqref{eq:ndcg}) is the SMS loss in Eq.~\eqref{eq:sms}, which weights both positive and negative pairs by their soft relevance. The size of the joint gain is consistent with the observation in~\cite{wray2021semantic} that graded-relevance objectives can reshape retrieval orderings well beyond what binary contrastive losses achieve.

Once the encoder has been fine-tuned end-to-end with SMS, however, further gains diminish quickly. TTA (\texttt{--flip --clip-length 32}) yields $+0.68$~nDCG with no training cost. Flip-averaging adapts the canonical recipe of~\cite{krizhevsky2012alexnet} to video and exploits the near-symmetry of kitchen scenes under horizontal reflection, while the longer clip length addresses the long tail of segment durations identified in Section~\ref{sec:materials:data}. Extending fine-tuning with six additional epochs on top of the initial 10-epoch run (16 epochs in total) adds a further $+0.20$~nDCG, suggesting the model is approaching the plateau of single-checkpoint training. This corresponds precisely to the regime in which model-averaging strategies tend to be most cost-effective and which motivates the future-work direction in Section~\ref{sec:future}. The practical implication of Table~\ref{tab:iterations} is that, under a tight compute budget, the bundle of encoder fine-tuning with relevance-aware supervision concentrates the bulk of the achievable improvement, while re-ranking and TTA are essentially free and should be applied by default, ensemble fusion of frozen baselines is also cheap but quickly hits a ceiling, and longer training is subject to rapidly diminishing returns. This ordering is consistent with what both the SMS-Loss authors~\cite{wang2024symmetricmultisimilaritylossepickitchens100} and the CR-CLIP authors~\cite{he2025crclip} report on the same benchmark, where most of the gap to the leaderboard is closed by the encoder and relevance-aware loss and subsequent refinements account for a smaller residual.
